\section*{अध्याय \ref{ch:g12:past-climates} --- अतीत की जलवायु पढ़ना}

\begin{solution}{exo:g12:past-climates:1}
(a) \omterm{def:g12:past-climates:proxy}{प्रॉक्सी}; (b) \omterm{def:g12:past-climates:ice-core}{बर्फ़ का कोर}; (c) \omterm{def:g12:past-climates:glacial-cycle}{अंतर्हिमनदीय काल}; (d) \omterm{def:g12:past-climates:pollen}{पराग-विश्लेषण};
(e) \omterm{def:g12:past-climates:lgm}{अंतिम हिमनदीय अधिकतम}।
\end{solution}

\begin{solution}{exo:g12:past-climates:2}
(a) ध्रुवीय बर्फ़: उस समय की वायु की कार्बन डाइऑक्साइड और मीथेन।
(b) झील के अवसाद या पीट: वनस्पति, यहाँ ऐसे पेड़ जिन्हें गरम
गर्मियाँ चाहिए। (c) समुद्री अवसाद: स्थल पर बर्फ़ का आयतन और \omterm{def:g10:layered-ocean:thermocline}{गहन
महासागर} का तापमान। (d) पेड़: वृद्धि की परिस्थितियाँ, वर्ष-दर-वर्ष।
\end{solution}

\begin{solution}{exo:g12:past-climates:3}
(a) \num{7000} वर्षों में $-50 - (-120) = \qty{70}{m}$: $\qty{70000}{mm}/\num{7000}
= \qty{10}{mm}$ प्रति वर्ष। (b) $10/3.7 \approx 2.7$: विहिमनदन की माध्य वृद्धि
2006--2018 की वृद्धि से लगभग 2.7 गुना तेज़ थी।
\end{solution}

\begin{solution}{exo:g12:past-climates:4}
एक शताब्दी दस दशक है: $\qty{2.0}{\degreeCelsius}$ प्रति शताब्दी। एक हज़ार
वर्ष सौ दशक हैं: $\qty{20}{\degreeCelsius}$ प्रति हज़ार वर्ष।
\end{solution}

\begin{solution}{exo:g12:past-climates:5}
(a) \num{23000} वर्ष: अक्ष का पुरस्सरण; \num{42000} वर्ष:
अक्ष का झुकाव; लगभग \num{100000} वर्ष: कक्षा की आकृति।
(b) एक मिलियन वर्षों में $\num{1000000}/\num{41000} \approx 24$ चक्र; $\num{1000000}/\num{100000}
= 10$ चक्र।
\end{solution}

\begin{solution}{exo:g12:past-climates:6}
(a) $266.6 + 7.83 \times (-9.4) \approx 266.6 - 73.6 = \qty{193}{ppm}$: मापे गए
\qty{190}{ppm} से केवल \qty{3}{ppm} ऊपर। (b) $266.6 + 7.83 \times 5.5
\approx \qty{310}{ppm}$, सबसे ऊँचे मापे गए मान से \qty{11}{ppm} ऊपर:
सबसे गरम बिंदु रेखा के नीचे पड़ता है। (c) $\Delta T = (427 - 266.6)/7.83 \approx
\qty{20}{\degreeCelsius}$। रेखा केवल यह बताती है कि दोनों राशियाँ प्राकृतिक चक्रों में,
आँकड़ों के परास के भीतर, साथ-साथ कैसे बदलीं
(\qty{-11}{\degreeCelsius} से \qty{+6}{\degreeCelsius}, 174 से
\qty{299}{ppm}); यह कोई नियम नहीं है, और उस परास से बहुत बाहर का मान
उससे नहीं पढ़ा जा सकता।
\end{solution}

\begin{solution}{exo:g12:past-climates:7}
(a) शाक, \qty{84}{\percent}: लगभग बिना पेड़ों का एक खुला, ठंडा भूदृश्य।
(b) ओक (\qty{42}{\percent}) और हेज़ल (\qty{34}{\percent}): गरम गर्मियों वाला
चौड़ी पत्ती का वन। (c) पहले भूर्ज फैलता है, फिर चीड़; हेज़ल और ओक
सबसे बाद में आते हैं। (d) पेड़ केवल उतनी तेज़ी से फैलते हैं जितनी तेज़ी से उनके बीज यात्रा करते हैं और उनके
नए पौधे बढ़ते हैं, उन स्थानों से जहाँ वे ठंड में बचे रहे; मृदा भी
बननी चाहिए। वन जलवायु से पीछे रह सकता है।
\end{solution}

\begin{solution}{exo:g12:past-climates:8}
(a) \num{8000} वर्षों में \qty{300}{cm}, यानी \qty{3000}{mm}: $3000/8000 =
\qty{0.375}{mm}$ प्रति वर्ष। (b) $\num{6000} + (470 - 300)/300 \times \num{8000}
\approx \num{6000} + \num{4530} \approx \num{10500}$ वर्ष। (c) $10/0.375
\approx 27$ वर्ष।
\end{solution}

\begin{solution}{exo:g12:past-climates:9}
(a) $7.16 - 4.11 = 3.05$ मिलियन वर्ग किलोमीटर, और $3.05/4.11 \approx
\qty{74}{\percent}$ अधिक भूमि। (b) हाँ: \omterm{def:g12:past-climates:lgm}{अंतिम हिमनदीय अधिकतम} पर समुद्र
आज से 125 से \qty{134}{m} नीचे था, जो मानचित्र पर बोर्नियो को एशिया से जोड़ने वाले
\qty{36}{m} से कहीं अधिक है। (c) मानचित्र आज के समुद्र-तल का उपयोग करता है, जो तब से
ठीक वैसा ही नहीं रहा (उस पर अवसाद निक्षेपित हुए हैं, और
भूमि थोड़ी खिसकी है); जाल भी मोटा है, हर 0.1
डिग्री पर एक बिंदु, और निम्नतम समुद्र-तल केवल लगभग दस मीटर की परिशुद्धता तक ज्ञात है।
\end{solution}

\begin{solution}{exo:g12:past-climates:10}
(a) \omterm{def:g12:past-climates:lgm}{अंतिम हिमनदीय अधिकतम} का मान, \qty{4.92}{\permil}, अधिक
है। यह अपेक्षित है: उस समय की बड़ी हिमचादरों ने ऑक्सीजन-18 में निर्धन
जल संचित किया, जिससे महासागर, और उसमें बढ़ी सीपियाँ, ऑक्सीजन-18 में
समृद्ध हो गईं; गहरा जल भी अधिक ठंडा था। (b) ताकि ठंडे, बर्फ़ीले समय
नीचे और गरम समय ऊपर पड़ें, जैसे किसी तापमान-वक्र पर।
\end{solution}

\begin{solution}{exo:g12:past-climates:11}
(a) सहसंबंध दिखाता है कि दोनों अभिलेख किसी सीधी रेखा के निकट
साथ-साथ उठते और गिरते हैं; वह यह नहीं दिखाता कि किसने किसे उत्पन्न किया, न यह कि दोनों में से किसी ने
किया। (b) गरम जलवायु ने वायु में कार्बन डाइऑक्साइड मुक्त की हो सकती है;
या कोई तीसरा कारण, जैसे पृथ्वी की कक्षा के परिवर्तन, दोनों को हिला सकता
था; या हर एक ने दूसरे को धकेला हो सकता है। (c) परिवर्तनों का समय,
कौन पहले आया, जो बर्फ़ और उसकी वायु के परिशुद्ध काल-निर्धारण से मिलता है; और यह
समझ कि कार्बन डाइऑक्साइड पृथ्वी के विकिरण पर कैसे काम करती है
(\cref{ch:g12:climate-change})।
\end{solution}

\begin{solution}{exo:g12:past-climates:12}
(a) 2.6 और 1.3 मिलियन वर्ष पहले के बीच $\num{1300000}/\num{41000} \approx 32$
चक्र; पिछले \num{700000} वर्षों में $\num{700000}/\num{98500} \approx 7$
चक्र। (b) \num{41000} वर्ष: अक्ष का झुकाव; लगभग
\num{100000} वर्ष: कक्षा की आकृति। (c) \num{41000} वर्ष के एक चक्र
का लगभग $\num{41000}/\num{2000} \approx 20$ बार नमूना लिया जाता है, जो उसे खींचने के लिए पर्याप्त है;
एक शताब्दी तक चला परिवर्तन दो नमूनों के बीच, या एक के भीतर, पड़ता है, और
लगभग कोई चिह्न नहीं छोड़ता।
\end{solution}

\begin{solution}{exo:g12:past-climates:13}
(a) $192/10 \approx 19$ और $198/10 \approx 20$: लगभग 20 गुना तेज़।
(b) \qty{4}{ppm} प्रति शताब्दी की तुलना में: $192/4 = 48$ से $198/4 \approx 50$
गुना; \qty{42}{ppm} की तुलना में: $192/42 \approx 4.6$ से $198/42 \approx 4.7$
गुना। आज की वृद्धि PETM की वृद्धि से लगभग 5 से 50 गुना तेज़ है।
(c) वृद्धि $2000 - 900 = \qty{1100}{ppm}$ है: PETM की सबसे तेज़ दर पर $1100/42 \approx 26$ शताब्दियाँ,
लगभग \num{2600} वर्ष; आज की दर पर $1100/195 \approx 5.6$
शताब्दियाँ, 600 वर्ष से कम।
\end{solution}

\begin{solution}{exo:g12:past-climates:14}
(a) 1 से \qty{1.5}{\degreeCelsius} का दस गुना: 10 से \qty{15}{\degreeCelsius}
प्रति हज़ार वर्ष, यानी 0.10 से \qty{0.15}{\degreeCelsius} प्रति दशक।
(b) आज की वैश्विक दर, \qty{0.20}{\degreeCelsius} प्रति दशक, ग्रीनलैंड के ऊपर के सबसे तेज़ तापन की 1.3 से 2
गुना है। (c) कोई एक स्थान पूरे विश्व के औसत से कहीं अधिक
बदल सकता है, जहाँ विभिन्न क्षेत्रों के परिवर्तन आंशिक रूप से
एक-दूसरे को काट देते हैं: वैश्विक दर की तुलना वैश्विक दर से की जाती है।
\end{solution}

\begin{solution}{exo:g12:past-climates:15}
(a) \qty{12}{m} प्रति हज़ार वर्ष पर \qty{14}{m} में $14/12 \approx
1.17$ हज़ार वर्ष लगते हैं, लगभग \num{1170} वर्ष: यानी लगभग \num{9800} वर्ष पहले।
(b) नहीं: बोर्नियो को एशिया से जोड़ने के लिए समुद्र को आज से कम से कम \qty{36}{m} नीचे
रहना था; \qty{-36}{m} पार करते ही स्थल-सेतु डूब गया, और
बोर्नियो फिर से द्वीप बन गया। (c) $36/12 = 3$ हज़ार वर्ष, लगभग
\num{3000} वर्ष; वास्तव में लगभग \num{8200} वर्ष पहले के बाद वृद्धि धीमी हो गई,
इसलिए अधिक समय लगा।
\end{solution}

\begin{solution}{pb:g12:past-climates:1}
\textbf{1.} स्वयं बर्फ़, अपने भारी समस्थानिकों के माध्यम से,
तापमान दर्ज करती है; उसके बुलबुलों की वायु कार्बन डाइऑक्साइड दर्ज करती है।

\textbf{2.} $5.5 - (-10.6) = \qty{16.1}{\degreeCelsius}$।

\textbf{3.} $299 - 174 = \qty{125}{ppm}$।

\textbf{4.} बिंदु एक सीधी रेखा के निकट पड़ते हैं: दोनों अभिलेख साथ-साथ उठे
और गिरे। यह नहीं बताता कि किसने किसे उत्पन्न किया।

\textbf{5.} अतीत के अभिलेख कभी सबसे पुराने समय को एक ओर, कभी दूसरी ओर रखकर खींचे जाते हैं;
उलटी ओर से पढ़ने पर तापन को शीतलन समझ लिया जाएगा।

\textbf{6.} \omterm{def:g12:past-climates:glacial-cycle}{हिमनदीय काल} किसी हिमयुग के भीतर एक ठंडा अंतराल है, जब
हिमचादरें अपने वर्तमान आकार से बहुत आगे तक बढ़ जाती हैं; \omterm{def:g12:past-climates:glacial-cycle}{अंतर्हिमनदीय काल} दो \omterm{def:g12:past-climates:glacial-cycle}{हिमनदीय कालों} के बीच
एक गरम अंतराल है, जिसमें हिमचादरें आज जितनी या उससे
छोटी होती हैं।

\textbf{7.} एक \omterm{def:g12:past-climates:glacial-cycle}{अंतर्हिमनदीय काल} में, जो लगभग \num{11000} वर्ष पहले शुरू हुआ।

\textbf{8.} $\num{800000}/\num{98000} \approx 8$: लगभग आठ चक्र।

\textbf{9.} कक्षा की आकृति।

\textbf{10.} \num{41000} वर्ष की लय, अक्ष के झुकाव का
आवर्तकाल।

\textbf{11.} $-0.4 - (-9.4) = \qty{9.0}{\degreeCelsius}$।

\textbf{12.} $264 - 190 = \qty{74}{ppm}$।

\textbf{13.} \qty{6}{\degreeCelsius}; अंटार्कटिक परिवर्तन, \qty{9.0}{\degreeCelsius},
बड़ा है, वैश्विक परिवर्तन का लगभग 1.5 गुना।

\textbf{14.} $\num{18000} - \num{11000} = \num{7000}$ वर्ष।

\textbf{15.} $6/7 \approx \qty{0.86}{\degreeCelsius}$ प्रति हज़ार वर्ष।

\textbf{16.} एक हज़ार वर्ष सौ दशक हैं: $0.86/100 \approx
\qty{0.0086}{\degreeCelsius}$ प्रति दशक।

\textbf{17.} $0.20 \times 100 = \qty{20}{\degreeCelsius}$ प्रति हज़ार वर्ष।

\textbf{18.} $20/0.86 \approx 23$ गुना।

\textbf{19.} $20/1.5 \approx 13$ और $20/1 = 20$: 13 से 20 गुना।

\textbf{20.} आज का वैश्विक तापन पिछले हिमयुग से बाहर निकलने वाले तापन से
\textbf{लगभग बीस गुना तेज़} है: उसकी माध्य दर का 23 गुना, उसकी सबसे तेज़ दर का
13 से 20 गुना।
\end{solution}
